% GENERATED FILE. Edit canonical lessons, then run npm run generate:books.
% canonical-source-hash: fnv1a64:57ffac36733c118c
% canonical-lessons: RU-C417-turniket, RU-C417-proyezdnoy, RU-C417-bezbiletnik, RU-C417-marshrut, RU-C417-marshrutka

\chapter{Metro and minibus}
\label{ch:ru-metro-and-minibus}
\hlchaptermodality{417}

\begin{chapteropening}
\textbf{By the end of this chapter:} \emph{I can get around town: the turnstile, a travel pass, a fare dodger, the route and the minibus.}

Complete the last lesson of chapter 417: I can get around town: the turnstile, a travel pass, a fare dodger, the route and the minibus.
\end{chapteropening}

\section[\texorpdfstring{\ru{турникет} (turnikét)}{турникет (turnikét)}]{\ru{турникет} (turnikét) — a turnstile}

\label{lesson:RU-C417-turniket}



\begin{quote}
Before the new one: say the Russian for a trolley, then the Russian for to present.
\end{quote}

\subsection*{You'll want to know: \ru{турникет}}
\textbf{\ru{турникет}} (\emph{turnikét}) — \textquotedblleft{}a turnstile\textquotedblright{}.

\begin{quote}
\ru{У} \ru{входа} \ru{в} \ru{метро} \ru{стоят} \ru{турникеты}.
\end{quote}

\begin{quote}
\emph{U vkhóda v metró stoyát turnikéty.}
\end{quote}

\begin{quote}
There are turnstiles at the metro entrance.
\end{quote}

Masculine.

\begin{cousinweb}
\emph{\ru{турникет}} was borrowed from French \emph{tourniquet} (EtymDB), built on \emph{tourner} 'to turn', which goes back to Latin \emph{tornāre} 'to turn on a lathe' (Lewis \& Short).
\end{cousinweb}

\begin{grammarlens}[title={What you've built}]
One more word for telling, explaining and describing at B1.
\end{grammarlens}

\subsection*{Your turn}
\begin{itemize}
\raggedright
  \item \emph{Say it:} \emph{turnikét}
  \item \emph{Say it:} \emph{turnikét}, once more
  \item \emph{Say it:} a sentence of your own with \emph{turnikét}
  \item \emph{Recall:} say the Russian for a trolley, then the Russian for to present, then say \emph{turnikét} again
\end{itemize}

\begin{tcolorbox}[breakable,colback=teal!4,colframe=teal!35!black,title={Before you move on}]
What does \ru{турникет} mean? (\textquotedblleft{}A turnstile\textquotedblright{}.) Say it once more.
\end{tcolorbox}
\section[\texorpdfstring{\ru{проездной} (proyezdnóy)}{проездной (proyezdnóy)}]{\ru{проездной} (proyezdnóy) — a travel pass, a season ticket}

\label{lesson:RU-C417-proyezdnoy}



\begin{quote}
Before the new one: say the Russian for to present, then the Russian for a turnstile.
\end{quote}

\subsection*{You'll want to know: \ru{проездной}}
\textbf{\ru{проездной}} (\emph{proyezdnóy}) — \textquotedblleft{}a travel pass, a season ticket\textquotedblright{}.

\begin{quote}
\ru{У} \ru{мамы} \ru{проездной} \ru{на} \ru{месяц}.
\end{quote}

\begin{quote}
\emph{U mámy proyezdnóy na mésyats.}
\end{quote}

\begin{quote}
Mum has a monthly travel pass.
\end{quote}

An adjective used as a noun, short for \emph{\ru{проездной} \ru{билет}}; it declines like an adjective.

\begin{cousinweb}
\emph{\ru{проездной}} is built on \emph{\ru{проезд}} 'a fare, a passage' with an adjective ending, from \emph{\ru{про}-} and the root of \emph{\ru{ездить}} 'to ride'. EtymDB lists \emph{\ru{проезд}} without a deeper source.
\end{cousinweb}

\begin{grammarlens}[title={What you've built}]
One more word for telling, explaining and describing at B1.
\end{grammarlens}

\subsection*{Your turn}
\begin{itemize}
\raggedright
  \item \emph{Say it:} \emph{proyezdnóy}
  \item \emph{Say it:} \emph{proyezdnóy}, once more
  \item \emph{Say it:} a sentence of your own with \emph{proyezdnóy}
  \item \emph{Recall:} say the Russian for to present, then the Russian for a turnstile, then say \emph{proyezdnóy} again
\end{itemize}

\begin{tcolorbox}[breakable,colback=teal!4,colframe=teal!35!black,title={Before you move on}]
What does \ru{проездной} mean? (\textquotedblleft{}A travel pass, a season ticket\textquotedblright{}.) Say it once more.
\end{tcolorbox}
\section[\texorpdfstring{\ru{безбилетник} (bezbilétnik)}{безбилетник (bezbilétnik)}]{\ru{безбилетник} (bezbilétnik) — a fare dodger}

\label{lesson:RU-C417-bezbiletnik}



\begin{quote}
Before the new one: say the Russian for a turnstile, then the Russian for a travel pass.
\end{quote}

\subsection*{You'll want to know: \ru{безбилетник}}
\textbf{\ru{безбилетник}} (\emph{bezbilétnik}) — \textquotedblleft{}a fare dodger\textquotedblright{}.

\begin{quote}
\ru{В} \ru{автобусе} \ru{поймали} \ru{двух} \ru{безбилетников}.
\end{quote}

\begin{quote}
\emph{V avtóbuse poymáli dvukh bezbilétnikov.}
\end{quote}

\begin{quote}
Two fare dodgers were caught on the bus.
\end{quote}

Built on \emph{\ru{без} \ru{билета}} — without a ticket.

\begin{cousinweb}
\emph{\ru{безбилетник}} is built from \emph{\ru{без}} 'without', \emph{\ru{билет}} 'a ticket' and \emph{-\ru{ник}} 'one who'. \emph{\ru{Билет}} came from French \emph{billet} 'a note, a ticket'; EtymDB does not list the compound.
\end{cousinweb}

\begin{grammarlens}[title={What you've built}]
One more word for telling, explaining and describing at B1.
\end{grammarlens}

\subsection*{Your turn}
\begin{itemize}
\raggedright
  \item \emph{Say it:} \emph{bezbilétnik}
  \item \emph{Say it:} \emph{bezbilétnik}, once more
  \item \emph{Say it:} a sentence of your own with \emph{bezbilétnik}
  \item \emph{Recall:} say the Russian for a turnstile, then the Russian for a travel pass, then say \emph{bezbilétnik} again
\end{itemize}

\begin{tcolorbox}[breakable,colback=teal!4,colframe=teal!35!black,title={Before you move on}]
What does \ru{безбилетник} mean? (\textquotedblleft{}A fare dodger\textquotedblright{}.) Say it once more.
\end{tcolorbox}
\section[\texorpdfstring{\ru{маршрут} (marshrút)}{маршрут (marshrút)}]{\ru{маршрут} (marshrút) — a route}

\label{lesson:RU-C417-marshrut}



\begin{quote}
Before the new one: say the Russian for a travel pass, then the Russian for a fare dodger.
\end{quote}

\subsection*{You'll want to know: \ru{маршрут}}
\textbf{\ru{маршрут}} (\emph{marshrút}) — \textquotedblleft{}a route\textquotedblright{}.

\begin{quote}
\ru{Этот} \ru{автобус} \ru{ходит} \ru{по} \ru{новому} \ru{маршруту}.
\end{quote}

\begin{quote}
\emph{Étot avtóbus khódit po nóvomu marshrútu.}
\end{quote}

\begin{quote}
This bus runs on a new route.
\end{quote}

Masculine. \emph{\ru{Маршрут}} is the planned way; \emph{\ru{дорога}} the road itself.

\begin{cousinweb}
\emph{\ru{маршрут}} was borrowed from German \emph{Marschroute} 'a line of march', built from \emph{Marsch} 'march' and \emph{Route} 'route' (EtymDB).
\end{cousinweb}

\begin{grammarlens}[title={What you've built}]
One more word for telling, explaining and describing at B1.
\end{grammarlens}

\subsection*{Your turn}
\begin{itemize}
\raggedright
  \item \emph{Say it:} \emph{marshrút}
  \item \emph{Say it:} \emph{marshrút}, once more
  \item \emph{Say it:} a sentence of your own with \emph{marshrút}
  \item \emph{Recall:} say the Russian for a travel pass, then the Russian for a fare dodger, then say \emph{marshrút} again
\end{itemize}

\begin{tcolorbox}[breakable,colback=teal!4,colframe=teal!35!black,title={Before you move on}]
What does \ru{маршрут} mean? (\textquotedblleft{}A route\textquotedblright{}.) Say it once more.
\end{tcolorbox}
\section[\texorpdfstring{\ru{маршрутка} (marshrútka)}{маршрутка (marshrútka)}]{\ru{маршрутка} (marshrútka) — a fixed-route minibus}

\label{lesson:RU-C417-marshrutka}



\begin{quote}
Before the new one: say the Russian for a fare dodger, then the Russian for a route.
\end{quote}

\subsection*{You'll want to know: \ru{маршрутка}}
\textbf{\ru{маршрутка}} (\emph{marshrútka}) — \textquotedblleft{}a fixed-route minibus\textquotedblright{}.

\begin{quote}
\ru{До} \ru{работы} \ru{я} \ru{езжу} \ru{на} \ru{маршрутке}.
\end{quote}

\begin{quote}
\emph{Do rabóty ya yézzhu na marshrútke.}
\end{quote}

\begin{quote}
I go to work by minibus.
\end{quote}

Colloquial for \emph{\ru{маршрутное} \ru{такси}}: a minibus that runs on a set route and stops on request.

\begin{cousinweb}
\emph{\ru{маршрутка}} is a colloquial short form of \emph{\ru{маршрутное} \ru{такси}} 'a route taxi', built on \emph{\ru{маршрут}} 'a route' (EtymDB lists \emph{\ru{маршрутка}}). \emph{\ru{Маршрут}} came from German \emph{Marschroute}.
\end{cousinweb}

\begin{grammarlens}[title={What you've built}]
One more word for telling, explaining and describing at B1.
\end{grammarlens}

\subsection*{Your turn}
\begin{itemize}
\raggedright
  \item \emph{Say it:} \emph{marshrútka}
  \item \emph{Say it:} \emph{marshrútka}, once more
  \item \emph{Say it:} a sentence of your own with \emph{marshrútka}
  \item \emph{Recall:} say the Russian for a fare dodger, then the Russian for a route, then say \emph{marshrútka} again
\end{itemize}

\begin{tcolorbox}[breakable,colback=teal!4,colframe=teal!35!black,title={Before you move on}]
What does \ru{маршрутка} mean? (\textquotedblleft{}A fixed-route minibus\textquotedblright{}.) Say it once more.
\end{tcolorbox}
